\documentclass[11pt]{article}
\usepackage[margin=1in]{geometry}
\usepackage[T1]{fontenc}
\usepackage[utf8]{inputenc}
\usepackage{amsmath,amssymb,amsthm}
\usepackage{enumitem}

\title{ICPC World Finals 2025\\G. Lava Moat}
\author{}
\date{}

\begin{document}
\maketitle

\section*{Problem Summary}

For a chosen elevation $c$, the lava moat must follow the contour line where the terrain has height $c$.
We must find the shortest such contour path that connects the western border edge of the rectangle to the
eastern border edge, or determine that no contour of any elevation can do so.

\section*{Key Observations}

\begin{itemize}[leftmargin=*]
    \item Fix an elevation $c$ that is not equal to any vertex height.
    In each triangle, the intersection with the horizontal plane $z=c$ is either empty or a single line
    segment joining the two triangle edges whose endpoints lie on opposite sides of the level.
    \item As $c$ moves between two consecutive vertex heights, the combinatorial structure of these contour
    segments does not change. Only their lengths change, and inside one triangle that length is a linear
    function of $c$.
    \item Therefore every connected contour component in such an interval has total length of the form
    $A c + B$.
    A minimum over a linear function on an interval is attained at an endpoint, so it is enough to inspect
    levels equal to vertex heights.
    \item A contour component is built by gluing triangle-edge crossings together.
    This suggests a DSU on triangulation edges rather than on vertices.
    When a triangle is active in a whole interval of elevations, it contributes one linear-length segment
    joining two crossing edges, so within that interval we should union those two edge components and add
    that segment length.
    \item A triangle is active on one or two disjoint intervals of the sorted vertex heights.
    These intervals can be inserted into a segment tree over the height order and processed with a rollback
    DSU.
\end{itemize}

\section*{Algorithm}

\begin{enumerate}[leftmargin=*]
    \item Sort all vertices by height. Each leaf of the segment tree corresponds to one vertex height.
    \item For every triangle, determine the height intervals in which the contour crosses that triangle.
    On each such interval, identify the two crossed triangle edges and compute the segment length in the
    form $a c + b$. Insert this ``active segment'' into the segment tree interval.
    \item Traverse the segment tree with a rollback DSU on triangulation edges.
    Each DSU component stores the sum of its currently active contour lengths in the form $A c + B$.
    \item At a leaf with elevation $c$:
    \begin{itemize}[leftmargin=*]
        \item if the west border edge and east border edge are already in the same DSU component, that
        component itself is a candidate moat;
        \item otherwise, the only new connection that can appear exactly at this height is through the leaf
        vertex.
        For every incident triangle we test whether the contour at level $c$ connects the vertex to the
        opposite triangle edge, compute that short segment length, and combine it with the best DSU
        component lengths reaching the west and east sides.
    \end{itemize}
    \item The minimum candidate over all leaves is the answer.
\end{enumerate}

\section*{Correctness Proof}

We prove that the algorithm returns the correct answer.

\paragraph{Lemma 1.}
Between two consecutive vertex heights, the set of crossed edges in every triangle is fixed, and every
active triangle contributes a contour segment whose length is linear in the elevation $c$.

\paragraph{Proof.}
Inside one triangle, the terrain is an affine function of $(x,y)$. Intersecting an affine plane with the
horizontal plane $z=c$ produces a line. As long as $c$ stays strictly between the same two vertex heights,
the line meets the same pair of triangle edges, so the contour segment endpoints move linearly along
those edges. Hence the segment length is linear in $c$. \qed

\paragraph{Lemma 2.}
For any open interval between consecutive vertex heights, the total length of every connected contour
component is a linear function of $c$ on that interval.

\paragraph{Proof.}
By Lemma 1, each active triangle contributes one linear-length segment throughout the interval.
The contour connectivity does not change inside the interval because no triangle changes its crossed edges.
Therefore a connected component is always the union of the same set of triangle segments, and its total
length is the sum of linear functions, hence linear. \qed

\paragraph{Lemma 3.}
Some optimal moat is attained at a vertex height.

\paragraph{Proof.}
Take any interval between consecutive vertex heights. By Lemma 2, every contour component length on
that interval is linear in $c$, so the minimum value of any fixed candidate path is attained at one of the
two endpoints of the interval. Since every feasible moat belongs to one of these candidates, an overall
optimum exists at a vertex height. \qed

\paragraph{Theorem.}
The algorithm outputs the correct answer for every valid input.

\paragraph{Proof.}
By Lemma 3 it is enough to inspect vertex heights.
For one such height $c$, the rollback DSU contains exactly the contour segments that remain active up to
that leaf, with exactly the correct length expressions, so every already-connected west-to-east contour
component is evaluated correctly.
If west and east are not already in the same component, the only new topological event at level $c$ can
occur at the vertex whose height equals $c$. The per-triangle checks performed at the leaf enumerate all
ways in which the contour can leave that vertex and join an already maintained DSU component.
Combining one such branch on the west side and one on the east side gives every possible moat passing
through that vertex, and no other new connection can appear at this exact height.
Therefore every feasible moat is considered and evaluated with its exact length, so the minimum printed
by the algorithm is the true optimum. \qed

\section*{Complexity Analysis}

Let $n$ be the number of vertices and $m$ the number of triangles.
Each triangle contributes $O(1)$ active intervals, each interval is inserted into the segment tree in
$O(\log n)$ nodes, and each inserted segment causes one rollback-DSU union.
Hence the total running time is $O((n+m)\log n)$, and the memory usage is $O((n+m)\log n)$ for the
segment tree plus $O(n+m)$ for the geometric data and the DSU.

\section*{Implementation Notes}

\begin{itemize}[leftmargin=*]
    \item The DSU is built on triangulation edges because contour segments connect edge crossings, not
    vertices.
    \item Rollback is essential: the same active contour segment belongs to many segment-tree nodes, so we
    need to undo unions when returning from recursion.
    \item All geometry is done with doubles; the required error tolerance is generous enough for this.
\end{itemize}

\end{document}
